% Problem Statement and Current-Law Gap. Built by the future pass from the
% instructions below. Do not process now. Names the gap and inventories current
% law without yet drafting statutory text (that is sections/statutory_text.tex).

[Write the problem statement as connected legislative prose in three movements,
synthesising the sources in the table below. Do not restate the table; reference
it. Cite laws by their bibliography keys so each statute resolves to its eCFR,
Cornell LII, or legislature URL in the rendered bibliography.]

[Movement 1, The problem. State plainly that Physical AI oncology trials now put
autonomous and semi-autonomous robots into direct contact with patients (surgical
humanoids, therapeutic positioning, diagnostic needle placement, rehabilitative
and companion systems), that the software governing these robot-patient
interactions is increasingly machine-generated, and that no current United States
law requires the verification of that code to clear before the code is generated
or executed. Frame this as the core safety-and-efficacy gap the bill closes. Draw
the embodiment-raises-the-stakes argument from the VVUQ-01 introduction (one body
concentrates all error potential) and the dual medical-product-versus-research-
tool classification problem from the national-platform regulatory landscape.]

[Movement 2, Current-law inventory and the gap. Walk the existing framework and
show, for each, where it stops short of a verification-before-generation mandate
for robot-patient interaction code. Cover the Federal Food, Drug, and Cosmetic
Act device authority (\texttt{usc-fdca}, \texttt{cfr-devices}), the investigational
pathways (\texttt{cfr-part312}, \texttt{cfr-part50}) and their Physical AI
adaptations (\texttt{kawchak2026cfr312}, \texttt{kawchak2026cfr50},
\texttt{kawchak2026iche6r3}), Medicare reasonableness and medical-necessity
determinations (\texttt{usc-1395y}), the Public Health Service Act
(\texttt{usc-phsa}), HIPAA privacy and security and electronic-records integrity
(\texttt{cfr-hipaa}, \texttt{cfr-part11}), and the federal nondiscrimination
guarantees (\texttt{usc-titlevi}, \texttt{usc-ada}). Note expressly that the
antidiscrimination provisions are codified at 42 U.S.C. (Title VI at
42 U.S.C. § 2000d; the Americans with Disabilities Act at 42 U.S.C. § 12101), so
the bill cites those titles. Use the national-platform government-framework and
regulatory-landscape sections for the three-branch authority map and the current
FDA AI posture.]

[Movement 3, The state patchwork and why uniform federal action is needed. Show
that several states have begun to regulate AI in health decisions but only at the
utilization-review and claims-handling layer, not at the robot-patient code-
verification layer: New York (\texttt{ny-a9149}), Texas (\texttt{tx-sb1822}),
California (\texttt{ca-sb1120}), and Florida (\texttt{fl-hb527}). Argue that this
patchwork leaves trial-conduct robot code unaddressed and motivates a federal
verification-before-generation standard that the states can then build on, with a
table reserved for the detailed crosswalk in \texttt{sections/prior\_law.tex}.]

% Model source-file table (parent paths at top, abbreviated leaves):
% \begin{table}[h]\centering
% \begin{tabularx}{\textwidth}{L{4.6cm} Y}
% \toprule
% \textbf{Source file (abbreviated leaf)} & \textbf{What to take from it} \\
% \midrule
% \multicolumn{2}{l}{\textit{Parent: papers/VVUQ-01/final-paper/sections/}} \\
% introduction.tex & Embodiment raises stakes; verification-first gap \\
% \multicolumn{2}{l}{\textit{Parent: physical-ai-oncology-trials/national-platform/new\_paper/final\_paper/sections/}} \\
% regulatory\_landscape.tex & Medical-product vs research-tool classification; gaps and pathways \\
% gov\_framework.tex & Three-branch authority map; current FDA AI framework \\
% \bottomrule
% \end{tabularx}
% \caption{Sources for the problem statement and current-law gap.}
% \end{table}
